\section{Conclusion}
\label{sec:conclusion}

Adherence to a national inpatient hyperglycemia guideline can be measured
automatically, at the level of its individual recommendations, across an entire
hospital network. Applying that measurement to 48,354 admissions in six Israeli
general hospitals shows that adherence improved after the guideline was issued,
that the improvement was modest and unevenly distributed across both
recommendations and centers, and that adherence is associated with lower 30-day
mortality and a lower burden of hyperglycemic complications once measured
baseline differences are balanced. The recommendations least well implemented
are largely the same ones whose adherence tracks subsequent outcomes: the
repeated insulin decisions distributed across a hospitalization,
rather than the single discrete actions taken on admission. That convergence,
together with the favorable shift in glycemic risk, is the evidence a guideline
needs to justify implementing it more completely. These are associations rather
than treatment effects, and they describe how a guideline was adopted rather
than how well any clinician practised. What the work establishes is that the
question is now answerable of an entire health system's records rather than of
a sample small enough to review by hand.
